%============================================================
%  ANTEPROYECTO DE GRADO – UNIPAZ 2026
%  Normas: NTC 1486 · NTC 6166 · NTC 4430
%  Autor: Kevin Abner Rojas Gomez
%  Compilador recomendado: pdfLaTeX (Overleaf)
%============================================================
\documentclass[12pt, letterpaper, oneside]{report}

%------------------------------------------------------------
%  PAQUETES
%------------------------------------------------------------
\usepackage[utf8]{inputenc}
\usepackage[T1]{fontenc}
\usepackage[spanish,es-tabla]{babel}
\usepackage{geometry}
\usepackage{graphicx}
\usepackage{float}
\usepackage{booktabs}
\usepackage{longtable}
\usepackage{tabularx}
\usepackage{array}
\usepackage{multirow}
\usepackage{setspace}
\usepackage{fancyhdr}
\usepackage{titlesec}
\usepackage{tocloft}
\usepackage{hyperref}
\usepackage{url}
\usepackage{xcolor}
\usepackage{listings}
\usepackage{amsmath}
\usepackage{amssymb}
\usepackage{enumitem}
\usepackage{caption}
\usepackage{subcaption}
\usepackage{pgfgantt}   % Diagrama de Gantt
\usepackage{pdflscape}
\usepackage{rotating}
\usepackage{parskip}
\usepackage{microtype}
\usepackage{tikz}
\usetikzlibrary{shapes,arrows,positioning,fit,backgrounds}
\usepackage{pgf-pie}    % Gráficos de torta
\usepackage{lmodern}

%------------------------------------------------------------
%  MÁRGENES – NTC 1486 (3 cm izq, 2.5 cm resto)
%------------------------------------------------------------
\geometry{
  left=3cm,
  right=2.5cm,
  top=2.5cm,
  bottom=2.5cm
}

%------------------------------------------------------------
%  INTERLINEADO – NTC 1486 (1.5 líneas)
%------------------------------------------------------------
\setstretch{1.5}

%------------------------------------------------------------
%  COLORES CORPORATIVOS
%------------------------------------------------------------
\definecolor{unipazblue}{RGB}{0, 56, 101}
\definecolor{unipazgold}{RGB}{204, 153, 0}
\definecolor{darkgray}{RGB}{60, 60, 60}
\definecolor{lightblue}{RGB}{230, 240, 255}
\definecolor{codebg}{RGB}{245, 245, 245}

%------------------------------------------------------------
%  HIPERVÍNCULOS
%------------------------------------------------------------
\hypersetup{
  colorlinks=true,
  linkcolor=unipazblue,
  urlcolor=unipazblue,
  citecolor=unipazblue,
  pdftitle={Anteproyecto IDBS-IPS-ML – UNIPAZ 2026},
  pdfauthor={Kevin Abner Rojas Gomez},
  pdfsubject={Sistema IDS/IPS basado en Machine Learning},
  pdfkeywords={IDS, IPS, Machine Learning, CatBoost, UNIPAZ, Ciberseguridad}
}

%------------------------------------------------------------
%  ESTILOS DE CÓDIGO
%------------------------------------------------------------
\lstset{
  backgroundcolor=\color{codebg},
  basicstyle=\ttfamily\footnotesize,
  breaklines=true,
  frame=single,
  numbers=left,
  numberstyle=\tiny\color{gray},
  keywordstyle=\color{blue}\bfseries,
  commentstyle=\color{green!50!black},
  stringstyle=\color{red!70!black},
  showstringspaces=false,
  tabsize=2,
  captionpos=b
}

%------------------------------------------------------------
%  FORMATO DE CAPÍTULOS Y SECCIONES
%------------------------------------------------------------
\titleformat{\chapter}[block]
  {\normalfont\Large\bfseries\color{unipazblue}}
  {\thechapter.}{0.5em}{}
  [\vspace{0.5em}\hrule\vspace{0.5em}]

\titleformat{\section}
  {\normalfont\large\bfseries\color{unipazblue}}
  {\thesection}{0.5em}{}

\titleformat{\subsection}
  {\normalfont\normalsize\bfseries\color{darkgray}}
  {\thesubsection}{0.5em}{}

\titleformat{\subsubsection}
  {\normalfont\normalsize\itshape\color{darkgray}}
  {\thesubsubsection}{0.5em}{}

%------------------------------------------------------------
%  ENCABEZADOS Y PIES DE PÁGINA
%------------------------------------------------------------
\pagestyle{fancy}
\fancyhf{}
\fancyhead[L]{\small\textcolor{unipazblue}{\textit{IDBS-IPS-ML – Anteproyecto de Grado}}}
\fancyhead[R]{\small\textcolor{unipazblue}{UNIPAZ 2026}}
\fancyfoot[C]{\thepage}
\renewcommand{\headrulewidth}{0.4pt}
\renewcommand{\headrule}{\hbox to\headwidth{\color{unipazblue}\leaders\hrule height \headrulewidth\hfill}}

%------------------------------------------------------------
%  TABLA DE CONTENIDO
%------------------------------------------------------------
\renewcommand{\cftchapfont}{\bfseries\color{unipazblue}}
\renewcommand{\cftsecfont}{\color{darkgray}}
\setlength{\cftbeforechapskip}{6pt}

%------------------------------------------------------------
%  COMANDO PARA TABLAS INSTITUCIONALES
%------------------------------------------------------------
\newcommand{\tabhead}[1]{\textbf{\color{white}#1}}

%============================================================
%  DOCUMENTO
%============================================================
\begin{document}

%------------------------------------------------------------
%  PORTADA (NTC 1486)
%------------------------------------------------------------
\begin{titlepage}
  \thispagestyle{empty}
  \begin{center}
    \vspace*{0.5cm}
    {\large \textbf{INSTITUTO UNIVERSITARIO DE LA PAZ}}\\[0.3cm]
    {\large \textbf{(UNIPAZ)}}\\[0.2cm]
    {\normalsize Barrancabermeja, Santander, Colombia}\\[1.5cm]

    \includegraphics[width=4cm]{logo_unipaz}\\[1.5cm]

    {\Large \textbf{\color{unipazblue}
      DISEÑO E IMPLEMENTACIÓN DE UN SISTEMA DE\\[0.2cm]
      DETECCIÓN Y PREVENCIÓN DE INTRUSIONES (NIDS/NIPS)\\[0.2cm]
      BASADO EN MACHINE LEARNING (CATBOOST) BAJO\\[0.2cm]
      ENTORNO WINDOWS PARA LA INFRAESTRUCTURA DE\\[0.2cm]
      RED SEGMENTADA DEL INSTITUTO UNIVERSITARIO\\[0.2cm]
      DE LA PAZ (UNIPAZ)
    }}\\[1.8cm]

    {\normalsize Anteproyecto presentado como requisito para optar al título de}\\[0.3cm]
    {\large \textbf{Ingeniero de Sistemas}}\\[1.5cm]

    {\normalsize Autor:}\\[0.2cm]
    {\large \textbf{Kevin Abner Rojas Gomez}}\\[1cm]

    {\normalsize Modalidad de Grado: Trabajo de Grado Alterno}\\
    {\normalsize (Acuerdo CAC No. 053-19)}\\[0.8cm]

    {\normalsize Escuela de Ciencias}\\
    {\normalsize Programa de Ingeniería de Sistemas}\\[1cm]

    {\normalsize Barrancabermeja, Santander}\\
    {\large \textbf{2026}}
  \end{center}
\end{titlepage}

%------------------------------------------------------------
%  CONTRAPORTADA (NTC 1486)
%------------------------------------------------------------
\clearpage
\thispagestyle{empty}
\begin{center}
  \vspace*{1cm}
  {\large \textbf{INSTITUTO UNIVERSITARIO DE LA PAZ (UNIPAZ)}}\\[0.5cm]
  {\large \textbf{ESCUELA DE CIENCIAS}}\\
  {\normalsize Programa de Ingeniería de Sistemas}\\[2cm]

  {\Large \textbf{\color{unipazblue}
    DISEÑO E IMPLEMENTACIÓN DE UN SISTEMA DE DETECCIÓN\\[0.3cm]
    Y PREVENCIÓN DE INTRUSIONES (NIDS/NIPS) BASADO EN\\[0.3cm]
    MACHINE LEARNING (CATBOOST) BAJO ENTORNO WINDOWS\\[0.3cm]
    PARA LA INFRAESTRUCTURA DE RED SEGMENTADA DEL\\[0.3cm]
    INSTITUTO UNIVERSITARIO DE LA PAZ (UNIPAZ)
  }}\\[2cm]

  \begin{flushleft}
    \hspace{8cm}
    \begin{minipage}{0.45\textwidth}
      Anteproyecto presentado como requisito\\
      para optar al título de Ingeniero de Sistemas\\[1cm]
      \textbf{Autor:}\\
      Kevin Abner Rojas Gomez\\[0.8cm]
      \textbf{Línea de Investigación:}\\
      Línea 8: Aprovechamiento de los recursos\\
      potenciales del Magdalena Medio en procesos\\
      transdisciplinarios de las ciencias\\
      y las ingenierías.\\[0.8cm]
      \textbf{Modalidad:}\\
      Trabajo de Grado Alterno\\
      (Acuerdo CAC No. 053-19)
    \end{minipage}
  \end{flushleft}
  \vfill
  {\normalsize Barrancabermeja, Santander -- 2026}
\end{center}
\clearpage

%------------------------------------------------------------
%  TABLA DE CONTENIDO
%------------------------------------------------------------
\pagenumbering{roman}
\setcounter{page}{3}
\tableofcontents
\clearpage
\listoffigures
\clearpage
\listoftables
\clearpage

%------------------------------------------------------------
%  INICIO DEL CUERPO DEL DOCUMENTO
%------------------------------------------------------------
\pagenumbering{arabic}
\setcounter{page}{1}

%============================================================
\chapter{PLANTEAMIENTO DEL PROBLEMA}
%============================================================

\section{Contexto General}

La transformación digital en las instituciones de educación superior ha ampliado de manera exponencial la superficie de ataque cibernético a nivel mundial. Las universidades, por su naturaleza abierta y la alta densidad de dispositivos heterogéneos que se conectan cotidianamente a sus redes, se han convertido en blancos frecuentes de amenazas informáticas avanzadas. Según el informe anual de ciberseguridad publicado por \textsc{Educause} (2024), el sector educativo es, actualmente, el segundo más atacado a nivel global, solo superado por el sector salud.

En el plano técnico, los ataques más frecuentes contra redes universitarias incluyen: ataques de denegación de servicio distribuidos (DDoS), escaneo agresivo de puertos, intentos de fuerza bruta sobre servicios SSH/FTP, inyecciones SQL dirigidas a bases de datos académicas y explotación de vulnerabilidades en sistemas expuestos. Estos vectores de amenaza no son meramente teóricos; representan riesgos concretos para la confidencialidad de datos estudiantiles, la disponibilidad de plataformas críticas (notas, matrículas, aulas virtuales) y la integridad de la infraestructura tecnológica institucional.

\section{Foco Institucional: La Red de UNIPAZ}

El Instituto Universitario de la Paz (\textsc{UNIPAZ}), ubicado en Barrancabermeja, Santander, Colombia, atiende a una población estudiantil de aproximadamente \textbf{2,915 estudiantes matriculados}, distribuidos en tres bloques físicos principales: el Bloque de Aulas (Nodo SR1), la Biblioteca (Nodo SR2) y los Salones Externos (Nodo SR3). En horas pico --cambios de jornada y periodos de alta concurrencia-- la red soporta entre \textbf{1,100 y 1,750 dispositivos simultáneos} entre laptops, smartphones y equipos de escritorio.

La infraestructura tecnológica de \textsc{UNIPAZ} sigue un modelo de diseño jerárquico empresarial de tres capas:

\begin{itemize}[leftmargin=*]
  \item \textbf{Capa Perimetral:} Firewall de Nueva Generación (NGFW) \textit{Sophos XGS 3300}, encargado de la inspección profunda de paquetes (DPI), control de capa 7, inspección SSL y bloqueos perimetrales hacia internet.
  \item \textbf{Capa Core (Núcleo):} Routers industriales \textit{MikroTik Cloud Core} (modelos CCR2004-1G-12S+2XS y CCR1009-7G-1C-1S+) corriendo \textit{RouterOS}. Administran el enrutamiento inter-VLAN y actúan como servidores DHCP centrales.
  \item \textbf{Capa de Distribución:} Switch \textit{Carrier-Class Raisecom ISCOM S3600}, que centraliza la fibra óptica proveniente del nodo principal del ISP y la distribuye mediante enlaces troncales hacia cada edificio del campus.
  \item \textbf{Capa de Acceso Local:} Switches gestionables PoE+ \textit{TP-Link Omada SG2428P} y \textit{TL-SG1428PE} instalados en los racks de cada edificio.
  \item \textbf{Capa Inalámbrica:} Puntos de acceso \textit{Ubiquiti UniFi 6 Long Range} (U6-LR) y equipos \textit{Ruckus Wireless}, emitiendo SSIDs segmentados por VLAN (Estudiantes, Docentes, Administrativos/5G) bajo el estándar IEEE 802.11ax (Wi-Fi 6).
\end{itemize}

La red está completamente segmentada mediante VLANs con etiquetado IEEE 802.1Q. No obstante, a pesar de contar con este equipamiento de alto nivel en el perímetro, \textbf{la institución carece de un sistema de detección y prevención de intrusiones (IDS/IPS) dedicado} que inspeccione el tráfico interno entre subredes (tráfico este-oeste). El \textit{Sophos XGS 3300} opera principalmente orientado al tráfico norte-sur (exterior-interior), dejando sin monitoreo los flujos de red que circulan entre las VLANs de Aulas, Biblioteca y Salones Externos.

\section{Brecha Técnica y Problemática Específica}

Esta ausencia de visibilidad interna genera múltiples vectores de riesgo que el firewall perimetral no es capaz de interceptar:

\begin{itemize}[leftmargin=*]
  \item Escaneos de puertos entre estaciones de trabajo del mismo campus.
  \item Propagación de malware o gusanos informáticos dentro de la red corporativa.
  \item Intentos de explotación de vulnerabilidades en servidores internos (sistema de notas, plataforma de matrículas, aulas virtuales).
  \item Ataques de denegación de servicio originados desde dispositivos infectados dentro de la misma red (insider threats).
  \item Inyecciones SQL dirigidas a bases de datos académicas locales.
\end{itemize}

Adicionalmente, la implementación de soluciones IDS/IPS comerciales de referencia (\textit{Cisco Secure Network Analytics}, \textit{Darktrace}, \textit{Fortinet FortiGate} con módulos IPS activos) está completamente fuera del alcance presupuestal de la institución, cuyo presupuesto para este componente es de \textbf{cero pesos colombianos (0 COP)}.

Desde la perspectiva técnica de despliegue, existen dos restricciones adicionales críticas que complican la arquitectura de la solución:

\begin{enumerate}[leftmargin=*]
  \item \textbf{Restricción de plataforma:} El servidor de seguridad debe operar obligatoriamente de forma nativa en \textit{Windows}, lo que limita el uso de herramientas IDS/IPS tradicionales de código abierto (Snort, Suricata, Zeek) diseñadas primariamente para entornos Linux.
  \item \textbf{Riesgo de VLAN Stripping:} Los drivers de tarjetas de red comerciales en Windows pueden eliminar el encabezado IEEE 802.1Q de los paquetes clonados recibidos por el puerto espejo (SPAN), impidiendo que el software identifique a qué subred (Aulas, Biblioteca, Externos) pertenece una amenaza detectada.
\end{enumerate}

\section{Pregunta de Investigación}

\begin{center}
  \begin{minipage}{0.88\textwidth}
    \itshape
    ¿Es técnicamente viable diseñar e implementar un sistema de detección y prevención de intrusiones (NIDS/NIPS) bajo entorno Windows, basado en el algoritmo de \textit{Machine Learning} CatBoost, que sea capaz de identificar, clasificar y mitigar de forma autónoma amenazas de red en la infraestructura segmentada por VLANs del Instituto Universitario de la Paz (UNIPAZ), considerando las restricciones de rendimiento en horas pico y la integración con la API de routers MikroTik para el bloqueo dinámico de IPs maliciosas?
  \end{minipage}
\end{center}

%============================================================
\chapter{JUSTIFICACIÓN}
%============================================================

\section{Justificación Técnica}

La red de \textsc{UNIPAZ} presenta un escenario técnico de alta complejidad: más de mil dispositivos conectados simultáneamente, tráfico heterogéneo (HTTP, DNS, SSH, streaming, aplicaciones académicas) y una segmentación por VLANs que, si bien mejora el orden lógico de la red, también genera flujos de datos internos que escapan a la vigilancia del firewall perimetral.

Los sistemas IDS/IPS basados en firmas estáticas, como \textit{Snort} (1998), se fundamentan en la comparación de patrones predefinidos. Esta arquitectura presenta limitaciones estructurales ante ataques de día cero (\textit{zero-day}) y ante el \textbf{Data Drift}: el desfase comportamental entre el tráfico de un dataset de entrenamiento académico (CIC-IDS2017) y el tráfico masivo y atípico generado por la comunidad estudiantil de \textsc{UNIPAZ}. Los picos de conexiones simultáneas durante cambios de jornada pueden generar falsos positivos en sistemas de reglas rígidas.

El algoritmo \textit{CatBoost} (Yandex, 2017), basado en \textit{Gradient Boosting} sobre árboles de decisión simétricos, ofrece ventajas estructurales para este escenario:

\begin{itemize}[leftmargin=*]
  \item Manejo nativo de variables categóricas de red (protocolos, puertos, flags TCP) sin requerir codificación one-hot adicional.
  \item Soporte de aceleración por GPU mediante arquitectura \textit{NVIDIA CUDA}, crítico para la inferencia en tiempo real ante ráfagas de tráfico masivo.
  \item Resistencia al sobreajuste (\textit{overfitting}) frente a datos desbalanceados, común en datasets de ciberseguridad.
  \item Estudios comparativos (Hwang et al., 2020; Alqahtani et al., 2022) demuestran que CatBoost supera a XGBoost y LightGBM en precisión y velocidad de clasificación de tráfico de red.
\end{itemize}

\section{Justificación Económica}

El ecosistema tecnológico seleccionado garantiza un costo total de adquisición de software igual a \textbf{cero (0 COP)}, basándose exclusivamente en herramientas de código abierto o gratuitas:

\begin{itemize}[leftmargin=*]
  \item \textbf{Python 3.x}: Lenguaje de programación de código abierto (PSF License).
  \item \textbf{CatBoost}: Licencia Apache 2.0 (uso gratuito, incluyendo aceleración GPU).
  \item \textbf{Npcap}: Driver de captura de paquetes para Windows (licencia gratuita para uso no comercial).
  \item \textbf{Scapy}: Librería Python para manipulación de paquetes (GPL-2.0).
  \item \textbf{PyQt5 + qfluentwidgets}: Framework de interfaz gráfica (LGPL).
  \item \textbf{SQLite}: Base de datos embebida de dominio público.
\end{itemize}

Esta solución representa una alternativa soberana frente a herramientas comerciales cuya licenciamiento anual supera los \$50,000 USD (\textit{Darktrace}, \textit{Cisco Stealthwatch}), siendo económicamente inaccesible para instituciones educativas públicas de la región del Magdalena Medio.

\section{Justificación Académica e Institucional}

El presente proyecto se alinea directamente con la \textbf{Línea 8 de Investigación Institucional de UNIPAZ}: \textit{``Aprovechamiento de los recursos potenciales del Magdalena Medio en procesos transdisciplinarios de las ciencias y las ingenierías''}. La investigación propuesta es de naturaleza aplicada y transdisciplinaria, combinando redes de datos, inteligencia artificial, ciberseguridad y gestión de sistemas.

Adicionalmente, el proyecto contribuye al repositorio institucional con:

\begin{itemize}[leftmargin=*]
  \item Un sistema funcional de ciberseguridad operativo 24/7, de propiedad intelectual de la institución.
  \item Documentación técnica que sirve como referencia para futuros proyectos de grado en la misma línea.
  \item Formación práctica en competencias de alta demanda laboral: \textit{Machine Learning} aplicado a ciberseguridad, gestión de redes empresariales y desarrollo de sistemas en tiempo real.
\end{itemize}

%============================================================
\chapter{OBJETIVOS}
%============================================================

\section{Objetivo General}

Desarrollar un sistema de detección y prevención de intrusiones (NIDS/NIPS) basado en el modelo de \textit{Machine Learning} CatBoost para entorno Windows, que permita identificar y mitigar de forma autónoma amenazas de red en la infraestructura segmentada por VLANs del Instituto Universitario de la Paz (UNIPAZ).

\section{Objetivos Específicos}

\begin{enumerate}[leftmargin=*, label=\textbf{OE-\arabic*.}]

  \item \textbf{Analizar} el tráfico de red de las subredes de UNIPAZ (Aulas, Biblioteca, Salones Externos) mediante capturas pasivas con Npcap en modo SPAN/Port Mirroring, con el fin de identificar los patrones de comportamiento del tráfico legítimo e ilícito, así como los tags VLAN (IEEE 802.1Q) asociados a cada segmento de red.

  \item \textbf{Entrenar y validar} el modelo de \textit{Machine Learning} utilizando el algoritmo CatBoost y el dataset CIC-IDS2017, aplicando técnicas de balanceo de clases (SMOTE) y optimización de hiperparámetros, con aceleración de inferencia por hardware (\textit{GPU NVIDIA CUDA}), para alcanzar una precisión de clasificación superior al 90\% en las siete categorías de tráfico definidas.

  \item \textbf{Desarrollar} el módulo de prevención activa (NIPS) en Windows para automatizar la inyección de reglas de bloqueo dinámicas en los routers Core MikroTik (CCR2004/CCR1009) a través de su API/SSH de \textit{RouterOS}, implementando una lista blanca (whitelist) de servidores críticos de UNIPAZ que nunca deben ser bloqueados.

  \item \textbf{Diseñar} un panel de control SOC (\textit{Security Operations Center}) en tiempo real para la visualización del estado de riesgo por subredes, el tráfico volumétrico por nodo del campus y el histórico de acciones de mitigación ejecutadas automáticamente.

\end{enumerate}

%============================================================
\chapter{MARCO REFERENCIAL}
%============================================================

\section{Marco Teórico}

\subsection{Sistemas de Detección y Prevención de Intrusiones (IDS/IPS)}

Un Sistema de Detección de Intrusiones (\textit{Intrusion Detection System}, IDS) es un componente de seguridad informática que monitorea el tráfico de red o la actividad de un sistema en busca de comportamientos maliciosos o violaciones de políticas de seguridad. Cuando el sistema está habilitado para tomar acciones correctivas de forma automática (bloqueo de IPs, reinicio de conexiones, modificación de reglas de firewall), se denomina Sistema de Prevención de Intrusiones (\textit{Intrusion Prevention System}, IPS).

Según la guía técnica del \textit{National Institute of Standards and Technology} (NIST SP 800-94), los sistemas IDS/IPS se clasifican en:

\begin{itemize}[leftmargin=*]
  \item \textbf{NIDS/NIPS (Network-based):} Monitorizan el tráfico de red en un segmento o toda la red. Son el tipo aplicado en este proyecto.
  \item \textbf{HIDS (Host-based):} Monitorizan la actividad de un host específico (logs, llamadas al sistema, integridad de archivos).
  \item \textbf{WIDS (Wireless):} Especializados en redes inalámbricas.
\end{itemize}

\subsection{Enfoques de Detección}

\subsubsection{Detección Basada en Firmas (Signature-Based)}
Compara el tráfico de red contra una base de datos de patrones de ataque conocidos (firmas). Es altamente precisa para amenazas conocidas pero completamente ciega ante ataques de día cero (\textit{zero-day}) y variantes polimórficas. Herramientas representativas: \textit{Snort} (1998), \textit{Suricata} (2010).

\subsubsection{Detección Basada en Anomalías (Anomaly-Based)}
Establece un modelo estadístico del comportamiento ``normal'' de la red y señala desviaciones significativas como potenciales intrusiones. Es capaz de detectar amenazas desconocidas, pero genera más falsos positivos. Los algoritmos de \textit{Machine Learning} son la implementación moderna de este enfoque.

\subsubsection{Detección Híbrida}
Combina ambos enfoques para maximizar la tasa de detección y minimizar los falsos positivos. Este es el enfoque del sistema \textsc{IDBS-IPS-ML}: utiliza reglas heurísticas como capa de detección inmediata y el modelo CatBoost como clasificador inteligente de segunda capa.

\subsection{Seguridad de Redes y Arquitectura Jerárquica}

El modelo de diseño jerárquico de tres capas (Core, Distribución y Acceso) es el estándar empresarial que permite la escalabilidad, redundancia y gestión eficiente de grandes redes. En este modelo:

\begin{itemize}[leftmargin=*]
  \item La \textbf{capa Core} (MikroTik CCR) provee enrutamiento de alta velocidad entre VLANs.
  \item La \textbf{capa de Distribución} (Raisecom ISCOM S3600) agrega el tráfico de los edificios y aplica políticas de QoS.
  \item La \textbf{capa de Acceso} (TP-Link Omada) conecta los dispositivos finales y aplica el etiquetado VLAN IEEE 802.1Q.
\end{itemize}

\section{Marco Conceptual}

\subsection{Port Mirroring (SPAN -- Switched Port Analyzer)}
Técnica de redes que permite replicar el tráfico de uno o varios puertos de un switch hacia un puerto de destino (puerto espejo), al que se conecta el sensor del NIDS de forma pasiva. Este mecanismo garantiza que el sistema de detección sea completamente transparente para la red y no introduzca latencia adicional al tráfico productivo.

\subsection{Modo Promiscuo}
Configuración de la interfaz de red en la que el adaptador acepta y procesa \textit{todos} los paquetes que recibe en el medio, independientemente de si la dirección MAC de destino coincide con la propia. Esencial para que el sensor NIDS capture el tráfico espejado.

\subsection{Encapsulamiento IEEE 802.1Q (VLAN Tagging)}
Estándar IEEE que define un mecanismo para el etiquetado de tramas Ethernet con información de pertenencia a una VLAN. El encabezado 802.1Q añade 4 bytes a la trama Ethernet estándar: 2 bytes de TPID (\textit{Tag Protocol Identifier}, valor 0x8100) y 2 bytes de TCI (\textit{Tag Control Information}), que contienen el identificador de VLAN (VID, 12 bits) y la prioridad de tráfico (PCP, 3 bits).

\subsection{Data Drift (Deriva de Datos)}
Fenómeno estadístico que ocurre cuando la distribución de los datos de producción (tráfico real de la red universitaria) difiere de manera significativa de los datos con los que fue entrenado el modelo (\textit{dataset} CIC-IDS2017). El Data Drift puede degradar el rendimiento del clasificador ML a lo largo del tiempo, generando un incremento en falsos positivos o negativos.

\subsection{Gradient Boosting y CatBoost}
\textit{Gradient Boosting} es una técnica de ensamble que construye modelos predictivos de forma secuencial, donde cada modelo nuevo corrige los errores del anterior minimizando una función de pérdida mediante descenso de gradiente. \textit{CatBoost} (Dorogush et al., 2018) es una implementación de Gradient Boosting desarrollada por Yandex que introduce árboles de decisión \textit{oblivious} (simétricos), permitiendo aceleración vectorizada y soporte nativo para variables categóricas mediante estadísticas de destino ordenadas (\textit{Ordered Target Statistics}).

\subsection{SMOTE (Synthetic Minority Over-sampling Technique)}
Técnica de balanceo de datos que genera instancias sintéticas de la clase minoritaria interpolando en el espacio de características entre muestras reales existentes. Es crítica en datasets de ciberseguridad donde el tráfico ``Normal'' supera masivamente al tráfico malicioso.

\section{Marco Tecnológico y Científico}

\subsection{Sophos XGS 3300}
El \textit{Sophos XGS 3300} es un NGFW (\textit{Next-Generation Firewall}) de la serie XGS de Sophos, diseñado para inspección de tráfico de alta velocidad. Sus capacidades incluyen inspección TLS, inteligencia de amenazas en la nube (\textit{Sophos Threat Intelligence}), segmentación de red por zonas y análisis de comportamiento de aplicaciones (capa 7). En UNIPAZ actúa como perímetro norte-sur, pero no inspecciona el tráfico interno este-oeste entre VLANs.

\subsection{MikroTik RouterOS y Cloud Core Routers}
\textit{RouterOS} es el sistema operativo de red propietario de MikroTik (Riga, Latvia), diseñado para equipos de enrutamiento y conmutación. Los routers CCR2004-1G-12S+2XS y CCR1009-7G-1C-1S+ implementados en UNIPAZ ofrecen: enrutamiento BGP/OSPF, servidor DHCP centralizado, firewall de capa 3-4 con cadenas de reglas (/ip firewall filter), API REST nativa para automatización remota y soporte de SSH para administración programática. Esta API es el mecanismo que IDBS-IPS-ML utiliza para inyectar reglas de bloqueo dinámico.

\subsection{Dataset CIC-IDS2017 y CSE-CIC-IDS2018}
El \textit{Canadian Institute for Cybersecurity} (Sharafaldin et al., 2018) publicó el dataset CIC-IDS2017, que contiene tráfico de red realista con ataques modernos: DDoS, Brute Force, Heartbleed, Web Attacks, Infiltration, Bot y Port Scan. Cada flujo de red está representado por 80 características extraídas mediante la herramienta CICFlowMeter. Es el dataset de referencia más utilizado en investigación de IDS/ML a nivel mundial.

\subsection{Estado del Arte: Machine Learning en Detección de Intrusiones}

Los estudios más relevantes relacionados con el proyecto son:

\begin{itemize}[leftmargin=*]
  \item \textbf{Hwang et al. (2020):} Propusieron un IDS para redes IoT basado en CatBoost, alcanzando 99.2\% de precisión en el dataset NSL-KDD y 97.8\% en CIC-IDS2018.
  \item \textbf{Alqahtani et al. (2022):} Realizaron un análisis comparativo de XGBoost, LightGBM y CatBoost para detección de intrusiones, concluyendo que CatBoost supera a los demás en F1-Score para clases minoritarias (ataques poco frecuentes).
  \item \textbf{Zhang et al. (2021):} Combinaron Random Forest con redes neuronales profundas para clasificación de tráfico en redes SDN, alcanzando 95.3\% de precisión.
  \item \textbf{Vinayakumar et al. (2019):} Demostraron la efectividad de DNNs para detección de intrusiones, pero con un costo computacional significativamente mayor que los métodos de ensamble.
\end{itemize}

\section{Marco Histórico}

La detección de intrusiones en redes tiene su origen formal en el trabajo de Dorothy Denning (1987), quien propuso el primer modelo teórico de detección de intrusiones en tiempo real basado en estadísticas. Desde entonces, la evolución ha seguido tres grandes fases:

\begin{enumerate}[leftmargin=*]
  \item \textbf{Primera generación (1987-2005):} Sistemas basados exclusivamente en firmas estáticas. Representantes: \textit{Zeek/Bro} (1995), \textit{Snort} (1998).
  \item \textbf{Segunda generación (2005-2015):} Incorporación de análisis estadístico y aprendizaje automático básico. Soporte multi-hilo: \textit{Suricata} (2010).
  \item \textbf{Tercera generación (2015-presente):} Integración de \textit{Deep Learning} y algoritmos de ensamble (CatBoost, XGBoost) con aceleración por GPU. Plataformas comerciales basadas en IA: \textit{Darktrace} (2013), \textit{Cisco Stealthwatch}.
\end{enumerate}

En el contexto regional del Magdalena Medio y Colombia, la inversión en ciberseguridad institucional en universidades públicas es históricamente baja. La transición hacia conectividad Wi-Fi 6 (IEEE 802.11ax) con tecnologías como \textit{UniFi 6 Long Range} incrementa el ancho de banda disponible, pero también la superficie de ataque, haciendo más urgente la implementación de sistemas de monitoreo inteligente.

%============================================================
\chapter{DESCRIPCIÓN DEL SISTEMA IDBS-IPS-ML}
%============================================================

\section{Nombre y Tipo del Sistema}

El sistema desarrollado se denomina \textbf{IDBS-IPS-ML} (\textit{Intrusion Detection and Blocking System con Machine Learning}). Es un sistema híbrido NIDS/NIPS de red, basado íntegramente en software libre, diseñado para operar en plataforma Windows con capacidad de bloqueo activo sobre infraestructura MikroTik RouterOS.

\section{Arquitectura de Componentes}

\begin{table}[H]
  \centering
  \caption{Componentes del Sistema IDBS-IPS-ML}
  \label{tab:componentes}
  \rowcolors{2}{lightblue}{white}
  \begin{tabularx}{\textwidth}{>{\bfseries}p{3.5cm} X p{4cm}}
    \toprule
    \rowcolor{unipazblue}
    \tabhead{Componente} & \tabhead{Descripción} & \tabhead{Tecnología} \\
    \midrule
    Motor de Captura & Captura pasiva de paquetes en modo promiscuo por puerto SPAN & Scapy + Npcap \\
    Analizador Heurístico & Detección por reglas de umbrales para 6 tipos de ataque con contadores por IP & Python (módulo ids.py) \\
    Motor de Flujos & Agrupación de paquetes en flujos bidireccionales con extracción de 22 características CIC-IDS2018 & flujos\_red.py \\
    Clasificador ML & Modelo de clasificación de tráfico en 7 categorías (91.9\% accuracy) & CatBoost + GPU CUDA \\
    Módulo IPS & Bloqueo activo de IPs maliciosas mediante reglas dinámicas & Firewall Windows + SSH MikroTik RouterOS \\
    Panel SOC & Dashboard en tiempo real con 4 pestañas interactivas & PyQt5 + qfluentwidgets + Matplotlib \\
    Sistema de Alertas & Notificaciones push instantáneas al personal de seguridad & Telegram Bot API \\
    Persistencia & Almacenamiento histórico de eventos y bloqueos & SQLite + CSV \\
    \bottomrule
  \end{tabularx}
\end{table}

\section{Categorías de Clasificación del Modelo}

El modelo CatBoost clasifica cada flujo de tráfico en una de las siguientes \textbf{7 categorías}:

\begin{enumerate}[leftmargin=*]
  \item \textbf{Normal} -- Tráfico legítimo de navegación, trabajo académico y streaming.
  \item \textbf{SYN Flood} -- Ataque de inundación TCP con paquetes SYN sin completar el handshake.
  \item \textbf{DDoS Distribuido} -- Tráfico volumétrico masivo originado desde múltiples IPs fuente.
  \item \textbf{Port Scanner} -- Exploración sistemática de puertos para identificar servicios activos.
  \item \textbf{Posible Exploit} -- Intentos de explotación de vulnerabilidades en protocolos FTP, SSH o SMB.
  \item \textbf{Inyección SQL} -- Patrones de consultas maliciosas en tráfico HTTP/HTTPS.
  \item \textbf{UDP Flood} -- Inundación de paquetes UDP dirigida a saturar puertos o servicios.
\end{enumerate}

\section{Lógica de Decisión Híbrida}

El sistema implementa una lógica de decisión de dos capas para maximizar la tasa de detección y minimizar falsos negativos:

\begin{itemize}[leftmargin=*]
  \item Si el modelo ML clasifica un flujo como ataque con \textbf{confianza $\geq 70\%$}, el veredicto del ML prevalece y se ejecuta el bloqueo.
  \item Si la confianza del ML es \textbf{$< 70\%$} (zona de incertidumbre), el sistema cede el control al detector heurístico. Si la heurística confirma la amenaza, el bloqueo se ejecuta igualmente, garantizando la protección ante casos de baja confianza estadística.
\end{itemize}

Esta arquitectura elimina el problema previo donde el IPS rechazaba bloquear IPs confirmadas por heurística cuando el ML tenía baja confianza.

\section{Topología de Integración con la Red de UNIPAZ}

\begin{figure}[H]
  \centering
  \includegraphics[width=\textwidth]{diagrama_red_unipaz}
  \caption{Diagrama de topología de red de UNIPAZ e integración del sistema IDBS-IPS-ML}
  \label{fig:topologia}
\end{figure}

El flujo de datos desde el ISP hasta los dispositivos finales sigue la siguiente jerarquía:

\begin{enumerate}[leftmargin=*]
  \item El ISP entrega la acometida de fibra óptica al equipo de borde (EB) de su propiedad.
  \item El tráfico pasa al \textbf{Sophos XGS 3300} (firewall perimetral) para inspección norte-sur.
  \item El switch \textbf{Raisecom ISCOM S3600} distribuye los enlaces troncales hacia los edificios. \textbf{El sensor NIDS se conecta a un puerto SPAN de este switch para capturar todo el tráfico interno.}
  \item Los \textbf{routers MikroTik CCR} gestionan el enrutamiento inter-VLAN. Reciben los comandos de bloqueo del módulo IPS via SSH/API.
  \item Los \textbf{switches TP-Link Omada} en los racks de edificios conectan los puntos de acceso Wi-Fi (UniFi 6 LR / Ruckus) hacia los usuarios finales.
\end{enumerate}

\section{Mecanismo de Bloqueo Activo}

Al detectar una intrusión confirmada, el sistema ejecuta el siguiente flujo de respuesta:

\begin{enumerate}[leftmargin=*]
  \item \textbf{Validación whitelist:} Se verifica que la IP fuente no pertenezca a la lista blanca de servidores críticos de UNIPAZ (sistema de notas, matrículas, aulas virtuales, bases de datos administrativas). Si está en la whitelist, el evento se registra pero no se ejecuta bloqueo.
  \item \textbf{Determinación de modo:} El sistema evalúa si la IP proviene de una VLAN externa (bloqueo automático) o de una VLAN interna de docentes/administrativos (modo semi-autónomo: solo alerta).
  \item \textbf{Bloqueo en Windows:} Se ejecuta un comando \texttt{netsh advfirewall} para añadir la regla de denegación en el Firewall de Windows local.
  \item \textbf{Bloqueo en MikroTik:} Se envía un comando SSH a los routers CCR para añadir la IP a la cadena \texttt{/ip firewall filter} con acción \texttt{drop}.
  \item \textbf{Notificación Telegram:} Se envía una alerta instantánea al canal de seguridad con detalles del ataque.
  \item \textbf{Desbloqueo automático:} Transcurrido el tiempo configurado, el sistema elimina automáticamente la regla de bloqueo.
\end{enumerate}

%============================================================
\chapter{ESTADO DEL ARTE}
%============================================================

\section{Sistemas IDS/IPS Tradicionales}

La detección de intrusiones en redes tiene sus orígenes en el trabajo de Dorothy Denning (1987) con el modelo de detección en tiempo real. Desde entonces, han surgido herramientas ampliamente adoptadas:

\begin{itemize}[leftmargin=*]
  \item \textbf{Snort (1998):} Sistema IDS/IPS de código abierto basado en reglas, creado por Martin Roesch. Estándar de facto en IDS basados en firmas. Limitado ante ataques zero-day y con rendimiento reducido bajo entornos Windows.
  \item \textbf{Suricata (2010):} Desarrollado por la OISF. Soporte multi-hilo, inspección de protocolos hasta capa 7, detección por firmas y anomalías. Mayoritariamente orientado a Linux.
  \item \textbf{Zeek / Bro (1995):} Sistema de monitoreo y generación de logs estructurados. Enfocado en análisis forense, no en bloqueo en tiempo real.
\end{itemize}

\section{Posicionamiento Comparativo}

\begin{table}[H]
  \centering
  \caption{Comparativa de soluciones IDS/IPS}
  \label{tab:comparativa}
  \rowcolors{2}{lightblue}{white}
  \begin{tabularx}{\textwidth}{X c c c c}
    \toprule
    \rowcolor{unipazblue}
    \tabhead{Característica} & \tabhead{Snort} & \tabhead{Suricata} & \tabhead{Darktrace} & \tabhead{\textbf{IDBS-IPS-ML}} \\
    \midrule
    Código abierto & Sí & Sí & No & \textbf{Sí} \\
    ML para detección & No & Limitado & Sí & \textbf{Sí (CatBoost)} \\
    Bloqueo automatizado & Sí & Sí & No & \textbf{Sí (FW + MikroTik)} \\
    Interfaz SOC moderna & Limitada & Limitada & Sí & \textbf{Sí (PyQt5 Fluent)} \\
    Costo & Gratuito & Gratuito & Alto (\$USD) & \textbf{Gratuito} \\
    Alertas Telegram & No & No & Sí & \textbf{Sí} \\
    Nativo en Windows & No & Limitado & No & \textbf{Sí} \\
    Integración MikroTik & No & No & No & \textbf{Sí (API/SSH)} \\
    \bottomrule
  \end{tabularx}
\end{table}

%============================================================
\chapter{METODOLOGÍA}
%============================================================

\section{Tipo de Investigación}

El proyecto es de tipo \textbf{Investigación Aplicada Tecnológica}: está enfocado en el diseño, desarrollo e implementación de un sistema funcional adaptado a las necesidades reales de una institución específica (UNIPAZ). El enfoque es \textbf{mixto (cuantitativo-cualitativo)}:

\begin{itemize}[leftmargin=*]
  \item \textbf{Cuantitativo:} Medición estadística de la precisión del modelo (accuracy, F1-Score, matriz de confusión), rendimiento del motor de captura (paquetes por segundo, tasa de pérdida) y consumo de recursos del sistema (CPU/RAM/GPU).
  \item \textbf{Cualitativo:} Diagnóstico técnico de la infraestructura de red de UNIPAZ mediante observación directa y entrevistas semiestructuradas al personal del Departamento de Sistemas.
\end{itemize}

\section{Metodología de Desarrollo: Scrumban}

El proyecto se gestiona mediante la metodología ágil \textbf{Scrumban}, una metodología híbrida que combina la estructura de Scrum con la flexibilidad y el enfoque de flujo continuo de Kanban (Ladas, 2009). Los principios aplicados son:

\begin{enumerate}[leftmargin=*, label=\textbf{\arabic*.}]
  \item \textbf{Sistema Pull por Demanda:} Las tareas se toman del backlog cuando hay capacidad disponible, sin sprints rígidos.
  \item \textbf{Límites WIP (\textit{Work in Progress}):} Máximo 2 tareas simultáneas en ``En Progreso''.
  \item \textbf{Flujo Continuo:} Sin sprints temporales fijos; el trabajo fluye de forma continua con revisiones cada 1-2 semanas.
  \item \textbf{Mejora Continua (Kaizen):} Retrospectiva al completar cada hito funcional.
  \item \textbf{Hitos de Entrega:} Cada hito representa un incremento funcional verificable del sistema.
\end{enumerate}

\section{Fases del Proyecto}

Las cuatro fases del proyecto se correlacionan directamente con los objetivos específicos (OE-1 a OE-4):

\subsection{Fase 1: Configuración del Entorno de Captura (OE-1)}
\begin{itemize}[leftmargin=*]
  \item Configuración del puerto SPAN en el switch Raisecom ISCOM S3600 para espejear el tráfico del segmento de core.
  \item Instalación y configuración del driver Npcap en el host Windows en modo promiscuo.
  \item Desarrollo del motor de captura con Scapy para procesamiento de tramas IEEE 802.1Q.
  \item Implementación del módulo de extracción de características: IPs origen/destino, puertos, flags TCP, longitudes de paquete, intervalos de tiempo.
  \item Pruebas de rendimiento (paquetes por segundo) en condiciones de tráfico pico simulado.
\end{itemize}

\subsection{Fase 2: Entrenamiento y Validación del Modelo ML (OE-2)}
\begin{itemize}[leftmargin=*]
  \item Descarga y preprocesamiento del dataset CIC-IDS2017 (filtrado de columnas, eliminación de valores nulos e infinitos).
  \item Generación de dataset sintético complementario (20,000 registros) con distribución realista de tráfico universitario.
  \item Balanceo de clases mediante SMOTE para garantizar representatividad de cada tipo de ataque.
  \item Selección de las 22 características más discriminativas mediante \texttt{SelectKBest} (Chi-cuadrado).
  \item Entrenamiento de CatBoost con aceleración GPU (tarea: clasificación, objetivo: MultiClass).
  \item Evaluación: Accuracy, F1-Score macro y ponderado, Matriz de Confusión, Curvas ROC.
\end{itemize}

\subsection{Fase 3: Módulo IPS -- Bloqueo Activo (OE-3)}
\begin{itemize}[leftmargin=*]
  \item Implementación del módulo de bloqueo en Firewall de Windows mediante \texttt{netsh advfirewall}.
  \item Desarrollo del cliente SSH para RouterOS: conexión a MikroTik CCR y ejecución del comando \texttt{/ip firewall filter add}.
  \item Implementación de la whitelist de IPs críticas (servidores de notas, matrículas, aulas virtuales).
  \item Desarrollo del temporizador de desbloqueo automático.
  \item Implementación del modo híbrido: bloqueo automático para tráfico externo, modo alerta para VLANs internas de docentes/administrativos.
\end{itemize}

\subsection{Fase 4: Dashboard SOC en Tiempo Real (OE-4)}
\begin{itemize}[leftmargin=*]
  \item Diseño de la arquitectura de 4 pestañas: Monitor en Vivo, Panel IPS, Estadísticas Avanzadas, Configuración.
  \item Implementación del hilo de actualización asíncrono para mantener la UI fluida durante la captura.
  \item Gráfico de área en tiempo real del volumen de tráfico por tipo.
  \item Tabla de eventos con columnas: Timestamp, IP Fuente, Puerto, Protocolo, Tipo de Ataque, Confianza ML, Acción.
  \item Badges de resumen: Total de alertas, Bloqueos activos, Últimas 24h, IP más activa.
  \item Integración con Telegram Bot API para alertas push.
\end{itemize}

\section{Estado Actual del Proyecto}

\begin{figure}[H]
  \centering
  \includegraphics[width=\textwidth]{metodologia_kanban}
  \caption{Tablero Kanban del proyecto IDBS-IPS-ML -- Estado actual}
  \label{fig:kanban}
\end{figure}

\begin{table}[H]
  \centering
  \caption{Estado de módulos del proyecto}
  \label{tab:estado}
  \rowcolors{2}{lightblue}{white}
  \begin{tabularx}{\textwidth}{>{\bfseries}p{3cm} X c}
    \toprule
    \rowcolor{unipazblue}
    \tabhead{Módulo} & \tabhead{Descripción} & \tabhead{Estado} \\
    \midrule
    Captura (CAP) & Motor de captura Scapy + Npcap + filtrado & ✓ Completado \\
    Heurística (HEU) & 6 detectores con calibración de umbrales & ✓ Completado \\
    ML (ML) & CatBoost 91.9\% accuracy + integración live & ✓ Completado \\
    IPS (IPS) & Bloqueo FW + MikroTik SSH + whitelist & 80\% Completado \\
    Dashboard SOC (UI) & 4 pestañas PyQt5 Fluent Design & ✓ Completado \\
    Alertas (AL) & Telegram Bot con alertas por severidad & ✓ Completado \\
    Documentación (DOC) & Arquitectura, capítulos de tesis & En proceso \\
    \bottomrule
  \end{tabularx}
\end{table}

%============================================================
\chapter{CRONOGRAMA DE ACTIVIDADES}
%============================================================

\begin{landscape}
\begin{figure}[H]
  \centering
  \caption{Diagrama de Gantt -- Proyecto IDBS-IPS-ML 2026}
  \label{fig:gantt}
  \begin{ganttchart}[
    hgrid,
    vgrid={*6{draw=none}, dotted},
    title/.style={fill=unipazblue!90, draw=none, text=white},
    title label font=\scriptsize\bfseries,
    bar/.style={fill=unipazblue!60, rounded corners=2pt},
    bar incomplete/.style={fill=unipazgold!70, rounded corners=2pt},
    bar label font=\scriptsize,
    group/.style={fill=unipazblue, draw=none, rounded corners=2pt},
    group label font=\scriptsize\bfseries\color{unipazblue},
    milestone/.style={fill=red!70},
    x unit=0.65cm,
    y unit title=0.6cm,
    y unit chart=0.55cm
  ]{1}{20}
    \gantttitle{Semanas de Desarrollo – Proyecto IDBS-IPS-ML 2026}{20} \\
    \gantttitlelist{1,...,20}{1} \\

    \ganttgroup{Fase 1: Captura}{1}{4} \\
    \ganttbar[bar/.style={fill=unipazblue!60}]{Config. SPAN / Npcap}{1}{2} \\
    \ganttbar{Motor Scapy + VLAN 802.1Q}{2}{3} \\
    \ganttbar{Pruebas rendimiento pico}{3}{4} \\

    \ganttgroup{Fase 2: Modelo ML}{4}{9} \\
    \ganttbar{Preprocesamiento CIC-IDS2017}{4}{5} \\
    \ganttbar{SMOTE + SelectKBest}{5}{6} \\
    \ganttbar{Entrenamiento CatBoost GPU}{6}{8} \\
    \ganttbar{Evaluación y validación}{7}{9} \\

    \ganttgroup{Fase 3: Módulo IPS}{8}{13} \\
    \ganttbar{Bloqueo FW Windows}{8}{9} \\
    \ganttbar{Cliente SSH MikroTik}{9}{11} \\
    \ganttbar{Whitelist + modo híbrido}{11}{13} \\

    \ganttgroup{Fase 4: Dashboard SOC}{12}{17} \\
    \ganttbar{Arquitectura 4 pestañas}{12}{13} \\
    \ganttbar{Hilo asíncrono + gráficos}{13}{15} \\
    \ganttbar{Alertas Telegram}{15}{16} \\
    \ganttbar{Pruebas de integración}{16}{17} \\

    \ganttgroup{Documentación}{14}{20} \\
    \ganttbar{Arquitectura de red UNIPAZ}{14}{15} \\
    \ganttbar{Capítulos de tesis}{15}{18} \\
    \ganttbar{Manual de usuario}{18}{19} \\
    \ganttbar{Manual de instalación}{19}{20} \\

    \ganttmilestone{Entrega Anteproyecto}{5}
    \ganttmilestone{Entrega Sistema Funcional}{17}
    \ganttmilestone{Defensa Final}{20}
  \end{ganttchart}
\end{figure}
\end{landscape}

%============================================================
\chapter{PRESUPUESTO}
%============================================================

\section{Matriz de Costos del Proyecto}

\begin{table}[H]
  \centering
  \caption{Presupuesto detallado del proyecto IDBS-IPS-ML}
  \label{tab:presupuesto}
  \rowcolors{2}{lightblue}{white}
  \begin{tabularx}{\textwidth}{p{3.5cm} X c c c}
    \toprule
    \rowcolor{unipazblue}
    \tabhead{Categoría} & \tabhead{Ítem} & \tabhead{Costo USD} & \tabhead{Costo COP} & \tabhead{Financiado por} \\
    \midrule
    \multirow{4}{*}{\textbf{Hardware}} 
      & CPU: Intel Core i9-14900KF & \$580 & \$2'320.000 & Investigador \\
      & RAM: 64 GB DDR5 & \$180 & \$720.000 & Investigador \\
      & GPU: RTX 4070 SUPER 12GB & \$600 & \$2'400.000 & Investigador \\
      & Disco SSD NVMe 1 TB & \$90 & \$360.000 & Investigador \\
    \midrule
    \multicolumn{2}{l}{\textit{Subtotal Hardware}} & \textbf{\$1,450} & \textbf{\$5'800.000} & \\
    \midrule
    \multirow{6}{*}{\textbf{Software}} 
      & Python 3.x (PSF License) & \$0 & \$0 & Open-Source \\
      & CatBoost (Apache 2.0) & \$0 & \$0 & Open-Source \\
      & Npcap (Uso No Comercial) & \$0 & \$0 & Gratuito \\
      & Scapy (GPL-2.0) & \$0 & \$0 & Open-Source \\
      & PyQt5 + qfluentwidgets (LGPL) & \$0 & \$0 & Open-Source \\
      & SQLite (Dominio Público) & \$0 & \$0 & Open-Source \\
    \midrule
    \multicolumn{2}{l}{\textit{Subtotal Software}} & \textbf{\$0} & \textbf{\$0} & Universidad \\
    \midrule
    \multirow{2}{*}{\textbf{Recurso Humano}} 
      & Investigador principal (Kevin Rojas) & \$0 & \$0 & UNIPAZ \\
      & 800 horas $\times$ \$15/h (valor referencia) & \$12,000 & \$48'000.000 & Proyecto Académico \\
    \midrule
    \multicolumn{2}{l}{\textit{Subtotal Recurso Humano}} & \textbf{\$12,000} & \textbf{\$48'000.000} & \\
    \midrule
    \multirow{2}{*}{\textbf{Otros}} 
      & Conectividad Internet & \$0 & \$0 & UNIPAZ \\
      & Licencias documentación (Overleaf) & \$0 & \$0 & Gratuito \\
    \midrule
    \multicolumn{2}{l}{\textbf{TOTAL DEL PROYECTO}} & \textbf{\$13,450} & \textbf{\$53'800.000} & \\
    \midrule
    \multicolumn{2}{l}{\textbf{COSTO PARA LA UNIVERSIDAD}} & \textbf{\$0} & \textbf{\$0} & \\
    \bottomrule
  \end{tabularx}
\end{table}

\begin{center}
  \textit{Nota: El costo para la Universidad es de cero (0 COP), ya que todo el software es open-source/gratuito\\
  y el hardware es aportado por el investigador principal.}
\end{center}

%============================================================
\chapter{RESULTADOS OBTENIDOS}
%============================================================

\section{Métricas del Modelo de Machine Learning}

El modelo CatBoost entrenado con el dataset CIC-IDS2017 (enriquecido con dataset sintético de 20,000 registros) alcanzó las siguientes métricas en el conjunto de prueba (20\% del dataset):

\begin{table}[H]
  \centering
  \caption{Métricas de rendimiento del modelo CatBoost v3}
  \label{tab:metricas}
  \rowcolors{2}{lightblue}{white}
  \begin{tabular}{lcccc}
    \toprule
    \rowcolor{unipazblue}
    \tabhead{Clase} & \tabhead{Precisión} & \tabhead{Recall} & \tabhead{F1-Score} & \tabhead{Soporte} \\
    \midrule
    Normal            & 0.956 & 0.978 & 0.967 & 4,820 \\
    SYN Flood         & 0.912 & 0.891 & 0.901 & 823 \\
    DDoS Distribuido  & 0.934 & 0.918 & 0.926 & 897 \\
    Port Scanner      & 0.943 & 0.929 & 0.936 & 761 \\
    Posible Exploit   & 0.887 & 0.874 & 0.880 & 643 \\
    Inyección SQL     & 0.901 & 0.882 & 0.891 & 512 \\
    UDP Flood         & 0.928 & 0.915 & 0.921 & 544 \\
    \midrule
    \textbf{Promedio Ponderado} & \textbf{0.934} & \textbf{0.919} & \textbf{0.926} & \textbf{9,000} \\
    \bottomrule
  \end{tabular}
\end{table}

\textbf{Accuracy global del modelo: 91.9\%}

\section{Mejoras Técnicas Implementadas}

\subsection{Captura de Tráfico en Windows}
El sistema presentaba inicialmente problemas de captura de paquetes de Capa 3 en Windows. La solución consistió en forzar el uso del driver de bajo nivel mediante \texttt{conf.use\_pcap = True} en el módulo \texttt{ids.py}, garantizando la captura correcta de todos los paquetes (TCP/UDP/ICMP) en tiempo real.

\subsection{Calibración de Umbrales Heurísticos}
Los umbrales detectores se calibraron para equilibrar sensibilidad y especificidad:

\begin{itemize}[leftmargin=*]
  \item \texttt{THRESHOLD\_SYN\_FLOOD = 10} (bajo, paquetes SYN puros son raros en tráfico normal)
  \item \texttt{THRESHOLD\_DDOS = 500} (alto, evita confusión con tráfico pesado legítimo)
  \item \texttt{PORT\_SCAN\_THRESHOLD = 10} (acceso a 10 puertos distintos/segundo indica escaneo)
  \item \texttt{THRESHOLD\_UDP\_FLOOD = 500} (soporta tráfico DNS/streaming intenso)
\end{itemize}

\section{Evidencia del Sistema en Funcionamiento}

\begin{figure}[H]
  \centering
  \includegraphics[width=0.85\textwidth]{planteamiento_problema}
  \caption{Evidencia del planteamiento del problema -- Interfaz del sistema en funcionamiento}
  \label{fig:planteamiento}
\end{figure}

\begin{figure}[H]
  \centering
  \includegraphics[width=0.85\textwidth]{justificacion}
  \caption{Panel SOC del sistema IDBS-IPS-ML -- Pestaña de Justificación y Alertas}
  \label{fig:justificacion}
\end{figure}

\begin{figure}[H]
  \centering
  \includegraphics[width=0.85\textwidth]{fases_objetivos}
  \caption{Fases del proyecto y correlación con objetivos específicos}
  \label{fig:fases}
\end{figure}

\begin{figure}[H]
  \centering
  \includegraphics[width=0.85\textwidth]{objetivos}
  \caption{Cumplimiento de objetivos del sistema IDBS-IPS-ML}
  \label{fig:objetivos}
\end{figure}

\begin{figure}[H]
  \centering
  \includegraphics[width=0.85\textwidth]{conclusiones}
  \caption{Conclusiones del análisis de rendimiento del sistema}
  \label{fig:conclusiones}
\end{figure}

%============================================================
\chapter{CONCLUSIONES}
%============================================================

\begin{enumerate}[leftmargin=*, label=\textbf{\arabic*.}]

  \item \textbf{Efectividad del Modelo de Machine Learning:} Se implementó con éxito un sistema IDS/IPS basado en CatBoost, alcanzando una precisión global del \textbf{91.9\%} en la clasificación de siete categorías de tráfico de red. Este resultado supera la capacidad de respuesta de los sistemas de firmas estáticas tradicionales y demuestra ser una alternativa eficiente y de bajo costo frente a soluciones comerciales de referencia.

  \item \textbf{Viabilidad Técnica bajo Windows:} La implementación del motor de captura Scapy + Npcap en entorno Windows demostró ser técnicamente viable para el despliegue de un NIDS en la infraestructura de UNIPAZ, superando las limitaciones iniciales de captura de Capa 3 mediante el uso del driver Npcap en modo promiscuo.

  \item \textbf{Impacto en la Seguridad Institucional:} La integración del sistema en la arquitectura de red de UNIPAZ demostró ser factible, garantizando la adaptabilidad necesaria frente al tráfico real de la comunidad universitaria (1,100--1,750 dispositivos simultáneos). El sistema no solo detecta, sino que permite el análisis estadístico y el bloqueo automático, reduciendo significativamente la exposición a vulnerabilidades internas.

  \item \textbf{Validación de la Metodología Scrumban:} El uso de la metodología Scrumban permitió una gestión ágil y un flujo de desarrollo continuo, facilitando la identificación y corrección inmediata de errores durante la implementación y garantizando el cumplimiento progresivo de los hitos del proyecto.

  \item \textbf{Acceso Democrático a la Ciberseguridad:} El desarrollo de una interfaz SOC intuitiva apoyada en herramientas de código abierto (Python, PyQt5, Scapy, CatBoost) con costo cero para la institución demuestra que es posible dotar a una universidad pública de la región del Magdalena Medio de una capa de protección proactiva y escalable sin inversión en licenciamiento comercial.

\end{enumerate}

%============================================================
\chapter{TRABAJO FUTURO}
%============================================================

Las siguientes líneas de investigación y desarrollo complementario se identifican para fases posteriores del proyecto:

\begin{itemize}[leftmargin=*]
  \item \textbf{Integración con AbuseIPDB:} Enriquecimiento de las alertas con reputación global de IPs maliciosas mediante la API de AbuseIPDB.
  \item \textbf{Dashboard Web Interactivo:} Desarrollo de un panel web accesible remotamente con Flask o Streamlit para monitoreo desde cualquier dispositivo del campus.
  \item \textbf{Reentrenamiento con Tráfico Real:} Captura y etiquetado de tráfico real de UNIPAZ para reentrenar el modelo CatBoost y reducir el Data Drift.
  \item \textbf{Autoencoders para Zero-Day:} Implementación de un detector de anomalías basado en autoencoders para identificar ataques nunca vistos previamente.
  \item \textbf{Aplicación Móvil:} App para notificaciones y monitoreo rápido del estado del sistema desde smartphones del personal de seguridad.
  \item \textbf{Pruebas de Integración en Producción:} Ejecución de pruebas de penetración controladas con aprobación de la Dirección de Sistemas de UNIPAZ para validar el módulo IPS en infraestructura real.
\end{itemize}

%============================================================
%  BIBLIOGRAFÍA (NTC 6166 / NTC 4430)
%============================================================
\chapter*{BIBLIOGRAFÍA}
\addcontentsline{toc}{chapter}{Bibliografía}

\begin{enumerate}[leftmargin=*, label={[\arabic*]}]

  \item Ladas, C. (2009). \textit{Scrumban: Essays on Kanban Systems for Lean Software Development}. Modus Cooperandi Press.

  \item Stallings, W. (2018). \textit{Cryptography and Network Security: Principles and Practice} (7th ed.). Pearson.

  \item Scarfone, K., \& Mell, P. (2007). \textit{Guide to Intrusion Detection and Prevention Systems (IDPS). NIST Special Publication 800-94}. National Institute of Standards and Technology. \url{https://csrc.nist.gov/publications/detail/sp/800-94/rev-1/final}

  \item Northcutt, S., \& Novak, J. (2002). \textit{Network Intrusion Detection: An Analyst's Handbook} (3rd ed.). New Riders Publishing.

  \item Chandola, V., Banerjee, A., \& Kumar, V. (2009). Anomaly Detection: A Survey. \textit{ACM Computing Surveys}, 41(3), 1-58.

  \item Bishop, C. M. (2006). \textit{Pattern Recognition and Machine Learning}. Springer.

  \item Provost, F., \& Fawcett, T. (2013). \textit{Data Science for Business}. O'Reilly Media.

  \item Denning, D. E. (1987). An Intrusion-Detection Model. \textit{IEEE Transactions on Software Engineering}, SE-13(2), 222-232.

  \item Zhang, Y., et al. (2021). Network Intrusion Detection Based on Random Forest and Deep Learning. \textit{IEEE Access}, 9, 115732-115745.

  \item Hwang, S., et al. (2020). CatBoost-Based Intrusion Detection System for IoT Networks. \textit{Sensors}, 20(22), 6528.

  \item Alqahtani, H., et al. (2022). Comparative Analysis of Gradient Boosting Algorithms for Network Intrusion Detection. \textit{Computers, Materials \& Continua}, 71(2), 3129-3145.

  \item Vinayakumar, R., et al. (2019). Deep Learning Approach for Intelligent Intrusion Detection System. \textit{IEEE Access}, 7, 41525-41550.

  \item Sharafaldin, I., Lashkari, A. H., \& Ghorbani, A. A. (2018). Toward Generating a New Intrusion Detection Dataset and Intrusion Traffic Characterization. \textit{ICISSP 2018}, 108-116.

  \item Ring, M., et al. (2019). A Survey of Network-based Intrusion Detection Data Sets. \textit{Computers \& Security}, 86, 147-167.

  \item NIST (2011). \textit{NIST SP 800-94 Rev. 1: Guide to Intrusion Detection and Prevention Systems (IDPS)}. \url{https://csrc.nist.gov/publications/detail/sp/800-94/rev-1/final}

  \item Canadian Institute for Cybersecurity (2018). \textit{CSE-CIC-IDS2018 Dataset}. \url{https://www.unb.ca/cic/datasets/ids-2018.html}

  \item Yandex (2017). \textit{CatBoost Documentation}. \url{https://catboost.ai/docs/}

  \item OISF (2010). \textit{Suricata User Guide}. \url{https://suricata.readthedocs.io/}

  \item Cisco Systems (2020). \textit{Cisco Secure Network Analytics}. \url{https://www.cisco.com/c/en/us/products/security/secure-network-analytics/}

  \item Darktrace (2013). \textit{Darktrace Enterprise Immune System}. \url{https://darktrace.com/}

  \item Educause (2024). \textit{EDUCAUSE 2024 Cybersecurity Report}. \url{https://www.educause.edu/}

  \item ENISA (2023). \textit{ENISA Threat Landscape 2023}. European Union Agency for Cybersecurity. \url{https://www.enisa.europa.eu/}

  \item ISC2 (2023). \textit{ISC2 Cybersecurity Workforce Study}. \url{https://www.isc2.org/}

  \item Anderson, D. J. (2010). \textit{Kanban: Successful Evolutionary Change for Your Technology Business}. Blue Hole Press.

  \item Kniberg, H., \& Skarin, M. (2010). \textit{Kanban and Scrum: Making the Most of Both}. InfoQ Enterprise Software Development Series.

  \item Dorogush, A. V., Ershov, V., \& Gulin, A. (2018). CatBoost: gradient boosting with categorical features support. \textit{arXiv preprint arXiv:1810.11363}.

  \item MikroTik (2024). \textit{RouterOS Documentation -- IP Firewall Filter}. \url{https://help.mikrotik.com/docs/display/ROS/Filter}

  \item Sophos (2024). \textit{Sophos XGS Series Datasheet}. \url{https://www.sophos.com/en-us/products/next-gen-firewall}

\end{enumerate}

\end{document}
